\section{Methods}

\subsection{Framework Overview and POMDP Formulation}
We model the patient-intervention ecosystem as a Partially Observable Markov Decision Process (POMDP). In psychiatric mHealth applications, the true latent clinical state (e.g., exact depression severity or momentary acute suicide risk) is inherently unobservable. Consequently, our RL agent operates strictly on an observable history, $\mathcal{H}_t$, which is restricted to the most recently recorded EMA values, the elapsed time since the last active patient response, and a cumulative measure of ``action fatigue'' (an index tracking the density of recent algorithmic interventions to prevent alert habituation).

\subsection{Continuous-Time Generative Modeling of Latent Trajectories}
To accurately capture the non-stationary, irregular, and episodic nature of psychiatric symptomatology, the unobserved latent clinical state, $X_t$, is generated using a continuous-time Neural Stochastic Differential Equation (Neural SDE):
$$dX_t = f_\theta(X_t, t)dt + g_\phi(X_t, t)dW_t$$
Within this formulation, the drift network $f_\theta$ governs the expected trajectory of symptom severity, while the diffusion network $g_\phi$ explicitly quantifies real-time clinical volatility. To simulate the abrupt and often non-linear decompensations characteristic of severe psychiatric disorders, this Neural SDE is hierarchically governed by a Hierarchical Dirichlet Process Hidden Markov Model (HDP-HMM). This structure induces sudden, non-parametric phase shifts between distinct latent regimes, effectively representing the transitions between stable ``Euthymic'' states and acute ``Crisis'' states without requiring a priori specification of the total number of clinical phenotypes. Following Li et al. (2020) and recent advancements in compact state-space modeling (arXiv:2508.23000, 2025), we constrain the SDE to ensure the predicted latent states remain strictly within valid, clinically interpretable domains (e.g., severity bounded between 0 and 1), preventing numerical divergence during volatile crisis phases. The implementation leverages Google Research's \texttt{torchsde} library for efficient, differentiable continuous-time integration.

\subsection{Informative Missingness and the ``Silent Smoke Detector'' Mechanism}
A critical structural challenge in EMA-based psychiatry is that non-response is profoundly Missing Not At Random (MNAR). As established by Lin et al. (2018), missing EMA data is an active behavioral phenotype; the absence of data often constitutes a high-acuity clinical signal (the ``Silent Smoke Detector'' effect). To mechanistically incorporate informative missingness into our framework, we model the temporal distribution of patient responses using a state-dependent Hawkes process. The baseline conditional intensity function of the observation process, $\lambda(t)$, is inversely coupled to the latent clinical state $X_t$ and the volatility measure derived from the Neural SDE. As latent severity or volatility increases (e.g., indicating a depressive collapse or manic agitation), $\lambda(t)$ dynamically decreases. This mathematically enforces a systemic reduction in patient engagement, realistically generating the characteristic I-MNAR patterns observed in high-risk psychiatric cohorts.

\subsection{The Dynamic Fidelity Controller and ESS-Constrained Optimization}
To preserve the OPE validity and scientific utility of the trial, the core RL agent is governed by a secondary, supervisory ``Dynamic Fidelity Controller.'' This controller continuously monitors the statistical power of the accumulating dataset by calculating Kish's Effective Sample Size (ESS) over a sliding temporal window. Rather than relying on trajectory-level ratios, we utilize Marginalized Importance Sampling (MIS) weights, $w(s, a)$, which directly estimate the stationary state-action density ratio between the RL agent's current behavioral policy and a predefined uniform baseline exploratory policy.

The real-time proxy for algorithmic fidelity and statistical utility is formulated as:
$$ESS_t = \frac{\left(\sum_{i=1}^n w(s_i, a_i)\right)^2}{\sum_{i=1}^n w(s_i, a_i)^2}$$
When the RL agent begins to ``game'' the environment by excessively exploiting specific, narrow state-action pairs, the variance of the MIS weights rapidly inflates, causing a sharp, detectable decline in $ESS_t$.

If $ESS_t$ falls below a critical threshold, $\tau_{crit}$, or if the underlying HDP-HMM infers a high probability of transition into an acute crisis state, the fidelity controller aggressively intervenes. It imposes a rigid constraint on policy optimization via a dynamically scaled Kullback-Leibler (KL) divergence penalty:
$$\mathcal{L}_{penalty} = \lambda_t D_{KL}\left(\pi_\theta(\cdot | \mathcal{H}_t) \,||\, \pi_{safe}(\cdot | \mathcal{H}_t)\right)$$
This penalty, akin to Reverse KL constraints utilized in Behavior Regularized Actor Critic (Wu et al., 2019), forces a rapid contraction of the agent's exploration boundaries. It immediately reverts the policy ($\pi_\theta$) to a conservative, pre-validated high-fidelity safety protocol ($\pi_{safe}$). Conversely, during confirmed euthymic periods combined with a robust ESS, the penalty coefficient $\lambda_t$ decays. This safely expands the exploration boundary, permitting the agent to harvest the essential stochastic counterfactual data required to power robust Marginalized Doubly Robust estimators at the trial's conclusion.
